\section{Review Method}
\label{sec:method}

\paragraph{Corpus.} The review is built around four core empirical studies published between 2023 and 2025. Each of them covers a different role of LLMs in the media context: simulation of human subjects~\cite{aher2023using}, representation of public opinion~\cite{santurkar2023whose}, generation of disinformation~\cite{vykopal2024disinformation} and detection of propaganda~\cite{gaeta2025towards}. Further sources are used for background and comparison: the SemEval-2020 Task~11 overview as the benchmark behind the detection study~\cite{dasanmartino2020semeval}, the MULTITuDE benchmark for machine-generated text detection~\cite{macko2023multitude}, conceptual work on information disorder and computational propaganda~\cite{wardle2017information, woolley2018computational, buchanan2021truth}, and the model papers that the studies build on~\cite{brown2020language, openai2023gpt4, ouyang2022training}.

\paragraph{Comparison dimensions.} Each core study was read with the same questions: What role does the LLM play? What task and data are used? Which models are evaluated? In which language? How is success measured? What is the main finding, and what limitation do the authors report? The answers form Table~\ref{tab:comparison} and the basis for the synthesis in Section~\ref{sec:comparative_synthesis}.

\paragraph{Verification of reported results.} All numbers in this paper were taken from the original publications and checked against the page or table where they are reported. Statements that could not be traced to a source were removed or weakened. For Gaeta et al.~\cite{gaeta2025towards} only the publisher's abstract was accessible; therefore no scores from that study are used. The plot in Figure~\ref{fig:detectors} only re-displays values from a published table. The only new data are those of the experiment in Section~\ref{sec:experiment}; its numbers are computed by the analysis script from the saved raw model outputs.
